\par\Needspace{6\baselineskip}\medskip
\hypertarget{native-130-article096}{}\pdfbookmark[2]{96.}{article-096}
\begin{center}96.\end{center}
The equations given in the preceding article are so elegant that they seem to leave nothing further to be desired. Nevertheless, some other formulas can be derived by which the orbital elements are determined still more elegantly and conveniently. The derivation of these formulas, however, is a little more involved.

We take up again the following equations from Art. 8, giving them new numbers for convenience.
\[
\mathrm{I.}\qquad \sin\tfrac12 v\mathpunct{.}\BareRoot{\dfrac{r}{a}}=\sin\tfrac12 E\mathpunct{.}\BareRoot{(1+e)}
\]
\SourceLeaf{131}
\[
\begin{aligned}
\mathrm{II.}\qquad \cos\tfrac12 v\mathpunct{.}\BareRoot{\dfrac{r}{a}}&=\cos\tfrac12 E\mathpunct{.}\BareRoot{(1-e)}\\[6pt]
\mathrm{III.}\qquad \sin\tfrac12 v'\mathpunct{.}\BareRoot{\dfrac{r'}{a}}&=\sin\tfrac12 E'\mathpunct{.}\BareRoot{(1+e)}\\[6pt]
\mathrm{IV.}\qquad \cos\tfrac12 v'\mathpunct{.}\BareRoot{\dfrac{r'}{a}}&=\cos\tfrac12 E'\mathpunct{.}\BareRoot{(1-e)}.
\end{aligned}
\]
We multiply I by $\sin\frac12(F+g)$ and II by $\cos\frac12(F+g)$. Adding the products gives
\[
\begin{aligned}
\cos\tfrac12(f+g)\mathpunct{.}\BareRoot{\dfrac{r}{a}}&=\sin\tfrac12 E\sin\tfrac12(F+g)\mathpunct{.}\BareRoot{(1+e)}\\[4pt]
&\quad+\cos\tfrac12 E\cos\tfrac12(F+g)\mathpunct{.}\BareRoot{(1-e)}
\end{aligned}
\]
or, since $\BareRoot{(1+e)}=\cos\frac12\varphi+\sin\frac12\varphi$ and $\BareRoot{(1-e)}=\cos\frac12\varphi-\sin\frac12\varphi$,
\[
\begin{aligned}
\cos\tfrac12(f+g)\mathpunct{.}\BareRoot{\dfrac{r}{a}}&=\cos\tfrac12\varphi\cos(\tfrac12 F-\tfrac12 G+g)\\[4pt]
&\quad-\sin\tfrac12\varphi\cos\tfrac12(F+G).
\end{aligned}
\]
In exactly the same way, multiplying III by $\sin\frac12(F-g)$ and IV by $\cos\frac12(F-g)$ and adding the products gives
\[
\begin{aligned}
\cos\tfrac12(f+g)\mathpunct{.}\BareRoot{\dfrac{r'}{a}}&=\cos\tfrac12\varphi\cos(\tfrac12 F-\tfrac12 G-g)\\[4pt]
&\quad-\sin\tfrac12\varphi\cos\tfrac12(F+G).
\end{aligned}
\]
Subtracting the preceding equation from this one gives
\[
\cos\tfrac12(f+g)\mathpunct{.}\left(\BareRoot{\dfrac{r'}{a}}-\BareRoot{\dfrac{r}{a}}\right)=2\cos\tfrac12\varphi\sin g\sin\tfrac12(F-G)
\]
or, on introducing the auxiliary angle $\omega$,
\[
\begin{aligned}
\mathrm{27)}\qquad &\cos\tfrac12(f+g)\operatorname{tang}2\omega\\[4pt]
&=\sin\tfrac12(F-G)\cos\tfrac12\varphi\sin g\mathpunct{.}\BareRoot[4]{\dfrac{aa}{rr'}}\mathpunct{\cdot}
\end{aligned}
\]
By exactly similar transformations, which we leave the experienced reader to work out, we find
\[
\begin{aligned}
\mathrm{28)}\qquad &\dfrac{\sin\frac12(f+g)}{\cos 2\omega}\\[4pt]
&=\cos\tfrac12(F-G)\cos\tfrac12\varphi\sin g\mathpunct{.}\BareRoot[4]{\dfrac{aa}{rr'}}\\[8pt]
\mathrm{29)}\qquad &\cos\tfrac12(f-g)\operatorname{tang}2\omega\\[4pt]
&=\sin\tfrac12(F+G)\sin\tfrac12\varphi\sin g\mathpunct{.}\BareRoot[4]{\dfrac{aa}{rr'}}\\[8pt]
\mathrm{30)}\qquad &\dfrac{\sin\frac12(f-g)}{\cos 2\omega}\\[4pt]
&=\cos\tfrac12(F+G)\sin\tfrac12\varphi\sin g\mathpunct{.}\BareRoot[4]{\dfrac{aa}{rr'}}\mathpunct{\cdot}
\end{aligned}
\]
Since the left-hand sides of these four equations are known quantities, from 27 and 28 we determine $\frac12(F-G)$ and $\cos\frac12\varphi\sin g\mathpunct{.}\BareRoot[4]{\dfrac{aa}{rr'}}=P$, and likewise from 29 and 30 we determine $\frac12(F+G)$ and $\sin\frac12\varphi\sin g\mathpunct{.}\BareRoot[4]{\dfrac{aa}{rr'}}=Q$. The ambiguity in \SourceLeaf{132}determining the angles $\frac12(F-G)$ and $\frac12(F+G)$ is to be resolved so that $P$ and $Q$ have the same sign as $\sin g$. From $P$ and $Q$ we then derive $\frac12\varphi$ and $\sin g\mathpunct{.}\BareRoot[4]{\dfrac{aa}{rr'}}=R$. From $R$ we can deduce $a=\dfrac{RR\BareRoot{rr'}}{\sin g^2}\mathpunct{\raisebox{.45ex}{,}}$ and $p=\dfrac{\sin f^2\BareRoot{rr'}}{RR}\mathpunct{\raisebox{.45ex}{,}}$ unless we prefer to use that quantity, which must come out $=\pm\BareRoot{\{2(l+\sin\frac12 g^2)\cos f\}}$, $=\pm\BareRoot{\{-2(L-\sin\frac12 g^2)\cos f\}}$, solely to check the calculation. In that case $a$ and $p$ are most conveniently determined by the formulas
\[
b=\dfrac{\sin f\mathpunct{.}\BareRoot{rr'}}{\sin g}\mathpunct{\raisebox{.45ex}{,}}\qquad a=\dfrac{b}{\cos\varphi}\mathpunct{\raisebox{.45ex}{,}}\qquad p=b\cos\varphi.
\]
Several equations from Arts. 88 and 95 may also be used at will to check the calculation. To these we add the following.
\[
\begin{aligned}
\dfrac{2\operatorname{tang}2\omega}{\cos 2\omega}\BareRoot{\dfrac{rr'}{aa}}&=e\sin G\sin g\\[6pt]
\dfrac{2\operatorname{tang}2\omega}{\cos 2\omega}\BareRoot{\dfrac{pp}{rr'}}&=e\sin F\sin f\\[6pt]
\dfrac{2\operatorname{tang}2\omega}{\cos 2\omega}&=\operatorname{tang}\varphi\sin G\sin f=\operatorname{tang}\varphi\sin F\sin g.
\end{aligned}
\]
Finally, the mean motion and the epoch of the mean anomaly are found in the same way as in the preceding article.
